% GENERATED FILE. Edit canonical lessons, then run npm run generate:books.
% canonical-source-hash: fnv1a64:075d7fc9415f57e4
% canonical-lessons: BN-C463-dushtu, BN-C463-rosik, BN-C463-bachal, BN-C463-thatta-kora, BN-C463-katukutu-deoya

\chapter{Mischief and Wit}
\label{ch:bn-mischief-and-wit}
\hlchaptermodality{463}

\begin{chapteropening}
\textbf{By the end of this chapter:} \emph{I can say that someone is naughty, witty or talkative.}

Complete the last lesson of chapter 463: I can say that someone is naughty, witty or talkative.
\end{chapteropening}

\section[\texorpdfstring{\bn{দুষ্টু} (dushṭu)}{দুষ্টু (dushṭu)}]{\bn{দুষ্টু} (dushṭu) — naughty; mischievous}

\label{lesson:BN-C463-dushtu}



\begin{quote}
Before the new one: say the Bengali for to shirk, then the Bengali for to dawdle.
\end{quote}

\subsection*{You'll want to know: \bn{দুষ্টু}}
\textbf{\bn{দুষ্টু}} — \emph{dushṭu} — \textquotedblleft{}naughty; mischievous\textquotedblright{}.

\begin{quote}
\bn{দুষ্টু} \bn{ছেলেটা} \bn{খাতায়} \bn{ছবি} \bn{আঁকল।}
\end{quote}

\begin{quote}
\emph{dushṭu chheleṭā khātāy chhôbi ā̃klô.}
\end{quote}

\begin{quote}
The naughty boy drew pictures in his exercise book.
\end{quote}

\bn{দুষ্টুমি}, mischief.

\begin{cousinweb}
\textbf{\bn{দুষ্টু}} (\emph{dushṭu}) comes from Sanskrit \emph{duṣṭa}, \textquotedblleft{}defiled, bad\textquotedblright{} (CDIAL 6472); Bengali softened it into a fond word for a mischievous child.
\end{cousinweb}

\begin{grammarlens}[title={What you've built}]
One more word for describing what you lived and felt.
\end{grammarlens}

\subsection*{Your turn}
\begin{itemize}
\raggedright
  \item \emph{Say it:} \emph{dushṭu}
  \item \emph{Say it:} \emph{dushṭu}, once more
  \item \emph{Say it:} \emph{dushṭu chheleṭā khātāy chhôbi ā̃klô.} — The naughty boy drew pictures in his exercise book.
  \item \emph{Recall:} say the Bengali for to shirk, then the Bengali for to dawdle, then say \emph{dushṭu} again
\end{itemize}

\begin{tcolorbox}[breakable,colback=teal!4,colframe=teal!35!black,title={Before you move on}]
What does \bn{দুষ্টু} mean? (\textquotedblleft{}Naughty; mischievous\textquotedblright{}.) Say it once more.
\end{tcolorbox}
\section[\texorpdfstring{\bn{রসিক} (rôsik)}{রসিক (rôsik)}]{\bn{রসিক} (rôsik) — witty; with a sense of humour}

\label{lesson:BN-C463-rosik}



\begin{quote}
Before the new one: say the Bengali for to dawdle, then the Bengali for naughty.
\end{quote}

\subsection*{You'll want to know: \bn{রসিক}}
\textbf{\bn{রসিক}} — \emph{rôsik} — \textquotedblleft{}witty; with a sense of humour\textquotedblright{}.

\begin{quote}
\bn{আমার} \bn{রসিক} \bn{মামার} \bn{গল্প} \bn{শুনে} \bn{আমরা} \bn{হাসলাম।}
\end{quote}

\begin{quote}
\emph{āmār rôsik māmār gôlpo shune āmrā hāslām.}
\end{quote}

\begin{quote}
We laughed at my witty uncle's stories.
\end{quote}

\begin{cousinweb}
\textbf{\bn{রসিক}} (\emph{rôsik}) is a tatsama, Sanskrit \emph{rasika}, \textquotedblleft{}tasteful, having a fine sense of taste\textquotedblright{} (Monier-Williams); \bn{রসিকতা} (taught), a joke, is its noun.
\end{cousinweb}

\begin{grammarlens}[title={What you've built}]
One more word for describing what you lived and felt.
\end{grammarlens}

\subsection*{Your turn}
\begin{itemize}
\raggedright
  \item \emph{Say it:} \emph{rôsik}
  \item \emph{Say it:} \emph{rôsik}, once more
  \item \emph{Say it:} \emph{āmār rôsik māmār gôlpo shune āmrā hāslām.} — We laughed at my witty uncle's stories.
  \item \emph{Recall:} say the Bengali for to dawdle, then the Bengali for naughty, then say \emph{rôsik} again
\end{itemize}

\begin{tcolorbox}[breakable,colback=teal!4,colframe=teal!35!black,title={Before you move on}]
What does \bn{রসিক} mean? (\textquotedblleft{}Witty; with a sense of humour\textquotedblright{}.) Say it once more.
\end{tcolorbox}
\section[\texorpdfstring{\bn{বাচাল} (bāchāl)}{বাচাল (bāchāl)}]{\bn{বাচাল} (bāchāl) — talkative; chatty}

\label{lesson:BN-C463-bachal}



\begin{quote}
Before the new one: say the Bengali for naughty, then the Bengali for witty.
\end{quote}

\subsection*{You'll want to know: \bn{বাচাল}}
\textbf{\bn{বাচাল}} — \emph{bāchāl} — \textquotedblleft{}talkative; chatty\textquotedblright{}.

\begin{quote}
\bn{বাচাল} \bn{মেয়েটা} \bn{সারাদিন} \bn{গল্প} \bn{করে।}
\end{quote}

\begin{quote}
\emph{bāchāl meyeṭā sārādin gôlpo kôre.}
\end{quote}

\begin{quote}
The chatty girl talks all day long.
\end{quote}

\begin{cousinweb}
\textbf{\bn{বাচাল}} (\emph{bāchāl}) is a tatsama, Sanskrit \emph{vācāla}, \textquotedblleft{}talkative, chattering\textquotedblright{}, said also of birds (Monier-Williams).
\end{cousinweb}

\begin{grammarlens}[title={What you've built}]
One more word for describing what you lived and felt.
\end{grammarlens}

\subsection*{Your turn}
\begin{itemize}
\raggedright
  \item \emph{Say it:} \emph{bāchāl}
  \item \emph{Say it:} \emph{bāchāl}, once more
  \item \emph{Say it:} \emph{bāchāl meyeṭā sārādin gôlpo kôre.} — The chatty girl talks all day long.
  \item \emph{Recall:} say the Bengali for naughty, then the Bengali for witty, then say \emph{bāchāl} again
\end{itemize}

\begin{tcolorbox}[breakable,colback=teal!4,colframe=teal!35!black,title={Before you move on}]
What does \bn{বাচাল} mean? (\textquotedblleft{}Talkative; chatty\textquotedblright{}.) Say it once more.
\end{tcolorbox}
\section[\texorpdfstring{\bn{ঠাট্টা} \bn{করা} (ṭhāṭṭā kôrā)}{ঠাট্টা করা (ṭhāṭṭā kôrā)}]{\bn{ঠাট্টা} \bn{করা} (ṭhāṭṭā kôrā) — to tease; to joke with}

\label{lesson:BN-C463-thatta-kora}



\begin{quote}
Before the new one: say the Bengali for witty, then the Bengali for talkative.
\end{quote}

\subsection*{You'll want to know: \bn{ঠাট্টা} \bn{করা}}
\textbf{\bn{ঠাট্টা} \bn{করা}} — \emph{ṭhāṭṭā kôrā} — \textquotedblleft{}to tease; to joke with\textquotedblright{}.

\begin{quote}
\bn{বন্ধুরা} \bn{অরুণের} \bn{নতুন} \bn{চুল} \bn{নিয়ে} \bn{ঠাট্টা} \bn{করল।}
\end{quote}

\begin{quote}
\emph{bôndhurā ôruner nôtun chul niye ṭhāṭṭā kôrlô.}
\end{quote}

\begin{quote}
His friends teased Arun about his new haircut.
\end{quote}

Past (he/she …): \bn{তিনি} \bn{নাতিকে} \bn{নিয়ে} \bn{ঠাট্টা} \bn{করলেন}, \emph{tini nātike niye ṭhāṭṭā kôrlen}.

\begin{cousinweb}
\textbf{\bn{ঠাট্টা} \bn{করা}} (\emph{ṭhāṭṭā kôrā}) is \emph{ṭhāṭṭā}, \textquotedblleft{}a joke\textquotedblright{}, which CDIAL lists for Bengali (CDIAL 5492), with \emph{kôrā}, \textquotedblleft{}to do\textquotedblright{} (CDIAL 2814).
\end{cousinweb}

\begin{grammarlens}[title={What you've built}]
One more word for describing what you lived and felt.
\end{grammarlens}

\subsection*{Your turn}
\begin{itemize}
\raggedright
  \item \emph{Say it:} \emph{ṭhāṭṭā kôrā}
  \item \emph{Say it:} \emph{ṭhāṭṭā kôrā}, once more
  \item \emph{Say it:} \emph{bôndhurā ôruner nôtun chul niye ṭhāṭṭā kôrlô.} — His friends teased Arun about his new haircut.
  \item \emph{Recall:} say the Bengali for witty, then the Bengali for talkative, then say \emph{ṭhāṭṭā kôrā} again
\end{itemize}

\begin{tcolorbox}[breakable,colback=teal!4,colframe=teal!35!black,title={Before you move on}]
What does \bn{ঠাট্টা} \bn{করা} mean? (\textquotedblleft{}To tease; to joke with\textquotedblright{}.) Say it once more.
\end{tcolorbox}
\section[\texorpdfstring{\bn{কাতুকুতু} \bn{দেওয়া} (kātukutu deoyā)}{কাতুকুতু দেওয়া (kātukutu deoyā)}]{\bn{কাতুকুতু} \bn{দেওয়া} (kātukutu deoyā) — to tickle}

\label{lesson:BN-C463-katukutu-deoya}



\begin{quote}
Before the new one: say the Bengali for talkative, then the Bengali for to tease.
\end{quote}

\subsection*{You'll want to know: \bn{কাতুকুতু} \bn{দেওয়া}}
\textbf{\bn{কাতুকুতু} \bn{দেওয়া}} — \emph{kātukutu deoyā} — \textquotedblleft{}to tickle\textquotedblright{}.

\begin{quote}
\bn{দুষ্টু} \bn{ভাইটা} \bn{আমাকে} \bn{কাতুকুতু} \bn{দিল।}
\end{quote}

\begin{quote}
\emph{dushṭu bhāiṭā āmāke kātukutu dilô.}
\end{quote}

\begin{quote}
My naughty brother tickled me.
\end{quote}

Past (he/she …): \bn{তিনি} \bn{নাতিকে} \bn{কাতুকুতু} \bn{দিলেন}, \emph{tini nātike kātukutu dilen}.

\begin{cousinweb}
\textbf{\bn{কাতুকুতু} \bn{দেওয়া}} (\emph{kātukutu deoyā}) is \emph{kātukutu}, an expressive word for tickling, with \emph{deoyā}, \textquotedblleft{}to give\textquotedblright{} (CDIAL 6141): to give someone a tickle.
\end{cousinweb}

\begin{grammarlens}[title={What you've built}]
One more word for describing what you lived and felt.
\end{grammarlens}

\subsection*{Your turn}
\begin{itemize}
\raggedright
  \item \emph{Say it:} \emph{kātukutu deoyā}
  \item \emph{Say it:} \emph{kātukutu deoyā}, once more
  \item \emph{Say it:} \emph{dushṭu bhāiṭā āmāke kātukutu dilô.} — My naughty brother tickled me.
  \item \emph{Recall:} say the Bengali for talkative, then the Bengali for to tease, then say \emph{kātukutu deoyā} again
\end{itemize}

\begin{tcolorbox}[breakable,colback=teal!4,colframe=teal!35!black,title={Before you move on}]
What does \bn{কাতুকুতু} \bn{দেওয়া} mean? (\textquotedblleft{}To tickle\textquotedblright{}.) Say it once more.
\end{tcolorbox}
